\section{Git and GitHub Foundations}

% =============================================================================
% GIT AND GITHUB FOUNDATIONS
% Insert directly into the existing Beamer project.
% No preamble, packages, theme definitions, or document-level commands included.
% =============================================================================

\begin{frame}{From Bash to Version Control}
\framesubtitle{The terminal operates files; Git manages their engineering history}
\footnotesize
\begin{center}
\colorbox{um6p-color-dark-green!18}{\parbox[c][0.92cm][c]{2.15cm}{\centering\textbf{Ubuntu}\\\scriptsize Completed}}
\quad$\Longrightarrow$\quad
\colorbox{um6p-color-dark-green!18}{\parbox[c][0.92cm][c]{2.15cm}{\centering\textbf{Terminal and Bash}\\\scriptsize Completed}}
\quad$\Longrightarrow$\quad
\colorbox{um6p-color-light-blue}{\parbox[c][1.02cm][c]{2.25cm}{\centering\textbf{Git and GitHub}\\\scriptsize Current foundation}}
\end{center}
\begin{block}{New Engineering Question}
How can several students modify the same robotics project while preserving its history, identifying each contribution, and keeping the GitHub repository reproducible?
\end{block}
\begin{alertblock}{Central Principle}
Git is not a replacement for careful engineering. It makes careful work visible, traceable, reviewable, and recoverable.
\end{alertblock}
\end{frame}

\begin{frame}{Why Robotics Needs Version Control}
\framesubtitle{Robotics systems evolve through many connected files and contributors}
\footnotesize
\begin{columns}[T,onlytextwidth]
\begin{column}{0.48\textwidth}
\begin{block}{Without Version Control}
\begin{itemize}\setlength{\itemsep}{0.2em}
\item files become \texttt{final}, \texttt{final2}, and \texttt{final\_fixed};
\item working configurations are overwritten;
\item changes cannot be attributed or reviewed;
\item two contributors may replace each other's work.
\end{itemize}
\end{block}
\end{column}
\begin{column}{0.48\textwidth}
\begin{exampleblock}{With Version Control}
\begin{itemize}\setlength{\itemsep}{0.2em}
\item each meaningful change becomes a recorded checkpoint;
\item experiments can be isolated on branches;
\item collaborators synchronize controlled changes;
\item earlier states remain available for investigation.
\end{itemize}
\end{exampleblock}
\end{column}
\end{columns}
\begin{alertblock}{Project Connection}
A change to a ROS 2 node may depend on a launch file, model, environment variable, or README instruction. Version control preserves these relationships as the system evolves.
\end{alertblock}
\end{frame}

\begin{frame}{What Is Version Control?}
\framesubtitle{A structured history of a project's meaningful states}
\footnotesize
\begin{center}
Working files $\longrightarrow$ Selected changes $\longrightarrow$ Recorded checkpoint $\longrightarrow$ Shared history
\end{center}
\begin{block}{Intuitive Definition}
Version control records how a collection of files changes over time. It helps engineers answer: \emph{What changed? Who changed it? Why was it changed? When was it changed?}
\end{block}
\pause
\begin{exampleblock}{Laboratory-Notebook Analogy}
A commit is similar to a signed experiment entry. It should describe one coherent change, preserve the evidence needed to understand it, and include a meaningful explanation.
\end{exampleblock}
\begin{alertblock}{Important Boundary}
Git records project history, but it does not automatically prove that code is correct. Tests, reviews, simulation evidence, and documentation remain necessary.
\end{alertblock}
\end{frame}

\begin{frame}{Git and GitHub Are Not the Same}
\framesubtitle{A local version-control tool and an online collaboration platform}
\footnotesize
\begin{columns}[T,onlytextwidth]
\begin{column}{0.47\textwidth}
\begin{block}{Git}
\begin{itemize}\setlength{\itemsep}{0.18em}
\item distributed version-control software;
\item runs on the student's computer;
\item records commits, branches, and merges;
\item can work locally without GitHub.
\end{itemize}
\end{block}
\end{column}
\begin{column}{0.47\textwidth}
\begin{exampleblock}{GitHub}
\begin{itemize}\setlength{\itemsep}{0.18em}
\item online hosting and collaboration platform;
\item stores shared Git repositories;
\item provides browsing, issues, reviews, and documentation;
\item connects contributors through remote repositories.
\end{itemize}
\end{exampleblock}
\end{column}
\end{columns}
\begin{alertblock}{One-Sentence Distinction}
Git manages the repository history; GitHub hosts and connects repositories for collaboration and delivery.
\end{alertblock}
\end{frame}

\begin{frame}{The Project Exists in Two Places}
\framesubtitle{The GitHub repository and the local clone have different roles}
\footnotesize
\begin{center}
\colorbox{um6p-color-light-blue}{\parbox[c][1.15cm][c]{3.0cm}{\centering\textbf{GitHub Remote}\\\scriptsize Shared project and collaboration point}}
\quad$\overset{\texttt{clone / pull}}{\underset{\texttt{push}}{\leftrightarrow}}$\quad
\colorbox{um6p-color-dark-green!18}{\parbox[c][1.15cm][c]{3.0cm}{\centering\textbf{Local Repository}\\\scriptsize Files, branches, commits, and local work}}
\end{center}
\begin{block}{Mental Model}
The remote is not a magical live folder. Students work in a local clone, create local commits, and then deliberately synchronize those commits with GitHub.
\end{block}
\begin{alertblock}{Practical Consequence}
Saving a file does not create a commit, and creating a commit does not upload it to GitHub. Each action is a separate engineering step.
\end{alertblock}
\end{frame}

\begin{frame}{Essential Git Vocabulary}
\framesubtitle{A project-centered mental model before using commands}
\scriptsize
\renewcommand{\arraystretch}{1.16}
\begin{tabularx}{\textwidth}{>{\bfseries\color{um6p-color-dark-blue}}p{0.18\textwidth} X}
\toprule
Repository & Project files together with the Git history and metadata. \\
Commit & A recorded snapshot of selected changes with author, time, and message. \\
Branch & A movable line of development used to isolate related work. \\
Remote & A named connection to another repository, commonly \texttt{origin} on GitHub. \\
Clone & A local copy of a repository, including its history and remote connection. \\
Pull & Retrieve remote changes and integrate them into the current local branch. \\
Push & Publish local commits to a remote branch. \\
Merge & Combine histories from different lines of development. \\
\texttt{.gitignore} & Rules describing untracked files that Git should normally ignore. \\
\bottomrule
\end{tabularx}
\end{frame}

\begin{frame}{Three Working Areas}
\framesubtitle{Git separates editing, selection, and permanent history}
\footnotesize
\begin{center}
\colorbox{black!7}{\parbox[c][1.0cm][c]{2.25cm}{\centering\textbf{Working Tree}\\\scriptsize Files being edited}}
$\xrightarrow{\texttt{git add}}$
\colorbox{um6p-color-light-blue}{\parbox[c][1.0cm][c]{2.25cm}{\centering\textbf{Staging Area}\\\scriptsize Selected next snapshot}}
$\xrightarrow{\texttt{git commit}}$
\colorbox{um6p-color-dark-green!18}{\parbox[c][1.0cm][c]{2.25cm}{\centering\textbf{Local History}\\\scriptsize Recorded commits}}
\end{center}
\begin{exampleblock}{Why Staging Exists}
A student may modify documentation, code, and temporary notes at the same time. Staging allows only the coherent project change to enter the next commit.
\end{exampleblock}
\begin{alertblock}{Key Distinction}
\texttt{git add} does not upload files. It selects content for the next local commit.
\end{alertblock}
\end{frame}

\begin{frame}[fragile]{Before Cloning: Verify the Destination}
\framesubtitle{Clone the repository once and into the intended parent directory}
\footnotesize
\begin{block}{Pre-Clone Checks}
\begin{verbatim}
cd ~
pwd
ls -la
which git
git --version
\end{verbatim}
\end{block}
\begin{exampleblock}{Question Before Execution}
Will \texttt{git clone} create \texttt{quard\_um6p\_intership} inside the current directory, or does a folder with that name already exist?
\end{exampleblock}
\begin{alertblock}{Common Mistake}
Do not clone the repository repeatedly inside itself or inside another accidental clone. If the target folder already exists, inspect it before taking action.
\end{alertblock}
\end{frame}

\begin{frame}{Access the Internship Repository on GitHub}
\framesubtitle{Understand the platform before copying commands}
\footnotesize
\begin{block}{Repository Address}
\centering
\url{https://github.com/abdtnourji/quard_um6p_intership}
\end{block}
\begin{columns}[T,onlytextwidth]
\begin{column}{0.48\textwidth}
\begin{exampleblock}{Inspect on GitHub}
\begin{itemize}\setlength{\itemsep}{0.18em}
\item repository name and owner;
\item default branch;
\item README and documentation;
\item folders and recent commits;
\item access permissions.
\end{itemize}
\end{exampleblock}
\end{column}
\begin{column}{0.48\textwidth}
\begin{block}{Obtain the URL}
Select \textbf{Code}, choose the documented authentication method, and copy the repository URL exactly. HTTPS is the simplest starting point when no SSH key workflow has been configured.
\end{block}
\end{column}
\end{columns}
\begin{alertblock}{Authenticity Check}
Confirm the repository owner and name before cloning. Similar names may refer to unrelated or outdated projects.
\end{alertblock}
\end{frame}

\begin{frame}[fragile]{Clone the Internship Repository}
\framesubtitle{Create one local working copy with history and remote configuration}
\footnotesize
\begin{block}{Purpose and Syntax}
\begin{verbatim}
git clone <repository-url>
\end{verbatim}
\end{block}
\begin{exampleblock}{Internship Command}
\begin{verbatim}
cd ~
git clone https://github.com/abdtnourji/quard_um6p_intership.git
cd quard_um6p_intership
\end{verbatim}
\end{exampleblock}
\begin{block}{What the Command Creates}
The project directory, working files, commit history, branches known at clone time, and a remote normally named \texttt{origin}.
\end{block}
\begin{alertblock}{If Authentication Is Requested}
Do not enter a GitHub account password as though it were a reusable Git password. Follow the repository's approved HTTPS credential or SSH-key procedure.
\end{alertblock}
\end{frame}

\begin{frame}[fragile]{Verify That the Clone Is Correct}
\framesubtitle{A successful-looking download still requires engineering checks}
\footnotesize
\begin{exampleblock}{Verification Sequence}
\begin{verbatim}
cd ~/quard_um6p_intership
pwd
ls -la
git status
git branch
git remote -v
git log --oneline -5
\end{verbatim}
\end{exampleblock}
\begin{block}{Expected Evidence}
The path identifies the intended repository, \texttt{.git} exists, Git recognizes a branch, the working tree is initially clean, \texttt{origin} points to the expected GitHub URL, and recent commits are visible.
\end{block}
\begin{alertblock}{Repository Test}
If \texttt{git status} reports ``not a git repository,'' verify the current directory and move to the cloned project root.
\end{alertblock}
\end{frame}

\begin{frame}[fragile]{Explore Before Modifying}
\framesubtitle{Understand the project structure and follow its documentation}
\footnotesize
\begin{columns}[T,onlytextwidth]
\begin{column}{0.48\textwidth}
\begin{block}{Filesystem View}
\begin{verbatim}
ls -la
find . -maxdepth 2 -type d | sort
find . -name "README*"
find . -name "package.xml"
\end{verbatim}
\end{block}
\end{column}
\begin{column}{0.48\textwidth}
\begin{exampleblock}{Git-Aware View}
\begin{verbatim}
git status
git branch
git log --oneline --graph -8
git remote -v
\end{verbatim}
\end{exampleblock}
\end{column}
\end{columns}
\begin{alertblock}{Professional Rule}
The README is part of the engineering interface. Read setup instructions, folder responsibilities, prerequisites, and known limitations before executing the pipeline.
\end{alertblock}
\end{frame}

\begin{frame}[fragile]{Use \texttt{git status} as the Dashboard}
\framesubtitle{Ask Git what it sees before and after every important action}
\footnotesize
\begin{block}{Purpose and Syntax}
\begin{verbatim}
git status
git status --short
\end{verbatim}
\end{block}
\begin{exampleblock}{Interpret the Categories}
\begin{itemize}\setlength{\itemsep}{0.18em}
\item \textbf{untracked}: Git has not been asked to include the file;
\item \textbf{modified}: tracked content differs from the last commit;
\item \textbf{staged}: content is selected for the next commit;
\item \textbf{clean}: no uncommitted change is detected.
\end{itemize}
\end{exampleblock}
\begin{alertblock}{Habit}
Run \texttt{git status} before editing, before staging, before committing, before pulling, and before pushing. It provides context, not just success or failure.
\end{alertblock}
\end{frame}

\begin{frame}[fragile]{Create a Small, Safe Change}
\framesubtitle{Practise the workflow with student-owned documentation}
\footnotesize
\begin{exampleblock}{Project Exercise}
\begin{verbatim}
cd ~/quard_um6p_intership
mkdir -p student_work
echo "Git workflow observations" > student_work/git_notes.txt
git status
\end{verbatim}
\end{exampleblock}
\begin{block}{Engineering Intention}
This creates a visible, low-risk change without modifying flight logic, simulation models, or shared configuration.
\end{block}
\begin{alertblock}{Repository Policy}
If the project defines a specific student folder, naming convention, or contribution process, follow that documentation instead of inventing a parallel structure.
\end{alertblock}
\end{frame}

\begin{frame}[fragile]{Inspect Changes with \texttt{git diff}}
\framesubtitle{Review content before selecting it for a commit}
\footnotesize
\begin{block}{Purpose and Syntax}
\begin{verbatim}
git diff              # unstaged tracked changes
git diff --staged     # changes selected for commit
\end{verbatim}
\end{block}
\begin{exampleblock}{Project Workflow}
\begin{verbatim}
echo "Repository inspected from the project root." \
  >> student_work/git_notes.txt
git diff
git status
\end{verbatim}
\end{exampleblock}
\begin{alertblock}{Important Detail}
An untracked file may not appear in ordinary \texttt{git diff}. Use \texttt{git status} to discover it, then stage it before reviewing with \texttt{git diff --staged}.
\end{alertblock}
\end{frame}

\begin{frame}[fragile]{Stage Only the Intended Change}
\framesubtitle{Build the next commit deliberately}
\footnotesize
\begin{block}{Purpose and Syntax}
\begin{verbatim}
git add <file-or-directory>
\end{verbatim}
\end{block}
\begin{exampleblock}{Preferred Beginner Workflow}
\begin{verbatim}
git add student_work/git_notes.txt
git status
git diff --staged
\end{verbatim}
\end{exampleblock}
\begin{alertblock}{Why Not Always Use \texttt{git add .}?}
Staging everything may include logs, generated files, credentials, unrelated edits, or large assets. Add explicit paths until you can explain every staged change.
\end{alertblock}
\end{frame}

\begin{frame}[fragile]{Create a Meaningful Commit}
\framesubtitle{Record one coherent engineering change with a useful message}
\footnotesize
\begin{block}{Purpose and Syntax}
\begin{verbatim}
git commit -m "concise description of the change"
\end{verbatim}
\end{block}
\begin{columns}[T,onlytextwidth]
\begin{column}{0.48\textwidth}
\begin{exampleblock}{Useful Message}
\begin{verbatim}
git commit -m \
  "docs: add student Git workflow notes"
\end{verbatim}
It states the area and result.
\end{exampleblock}
\end{column}
\begin{column}{0.48\textwidth}
\begin{alertblock}{Weak Messages}
\texttt{update}, \texttt{changes}, \texttt{final}, and \texttt{fix stuff} do not explain engineering intent.
\end{alertblock}
\end{column}
\end{columns}
\begin{block}{After Committing}
Run \texttt{git status} and \texttt{git log --oneline -5}. The commit exists locally until it is pushed.
\end{block}
\end{frame}

\begin{frame}[fragile]{Read the Project History}
\framesubtitle{Commits form an engineering narrative, not merely a backup list}
\footnotesize
\begin{block}{Purpose and Useful Views}
\begin{verbatim}
git log
git log --oneline
git log --oneline --graph --decorate --all
\end{verbatim}
\end{block}
\begin{exampleblock}{Questions the History Can Answer}
\begin{itemize}\setlength{\itemsep}{0.18em}
\item Which change introduced a launch instruction?
\item When was a model or configuration updated?
\item Which branch contains a student's contribution?
\item What commit preceded a newly observed problem?
\end{itemize}
\end{exampleblock}
\begin{alertblock}{Quality Signal}
Small coherent commits with descriptive messages are easier to review, merge, test, and recover than one large unexplained commit.
\end{alertblock}
\end{frame}

\begin{frame}[fragile]{Branches Isolate Lines of Work}
\framesubtitle{Experiment without immediately changing the shared main line}
\footnotesize
\begin{block}{Inspect and Create}
\begin{verbatim}
git branch
git switch -c docs/student-git-notes
\end{verbatim}
\texttt{git branch} lists branches. \texttt{git switch -c} creates and enters a new branch.
\end{block}
\begin{exampleblock}{Move Between Existing Branches}
\begin{verbatim}
git switch main
git switch docs/student-git-notes
\end{verbatim}
\end{exampleblock}
\begin{alertblock}{Before Switching}
Check \texttt{git status}. Uncommitted changes may follow you, conflict with the destination, or prevent the switch. Branch names should describe the task, not the contributor's mood.
\end{alertblock}
\end{frame}

\begin{frame}{What Does Merge Mean?}
\framesubtitle{Combine compatible histories after review and verification}
\footnotesize
\begin{center}
\texttt{main}: A $\longrightarrow$ B $\longrightarrow$ C

\vspace{0.15cm}
\hspace{1.0cm}$\searrow$ D $\longrightarrow$ E \quad \texttt{feature branch}

\vspace{0.15cm}
$\Downarrow$

\texttt{main}: A $\longrightarrow$ B $\longrightarrow$ C $\longrightarrow$ M
\end{center}
\begin{block}{Intuition}
A merge combines changes from different lines of development. If both lines modify incompatible parts of the same content, Git asks an engineer to resolve the conflict.
\end{block}
\begin{alertblock}{Responsibility}
A conflict is not Git failing. It is Git refusing to guess which engineering intention is correct. Resolve it by understanding both changes, then test the combined result.
\end{alertblock}
\end{frame}

\begin{frame}[fragile]{Inspect the Remote Connection}
\framesubtitle{Know where pulls come from and where pushes go}
\footnotesize
\begin{block}{Purpose and Syntax}
\begin{verbatim}
git remote -v
\end{verbatim}
\end{block}
\begin{exampleblock}{Expected Interpretation}
\texttt{origin} is the conventional name created by cloning. The output normally shows fetch and push URLs for the GitHub repository.
\end{exampleblock}
\begin{alertblock}{Wrong-Remote Prevention}
Before your first push, confirm that the owner and repository name are correct and that you have authorized contribution access. Do not replace \texttt{origin} randomly to bypass a permission problem.
\end{alertblock}
\end{frame}

\begin{frame}[fragile]{Synchronize Before Publishing}
\framesubtitle{Bring in remote work before sending your own commits}
\footnotesize
\begin{block}{Safe Beginner Sequence}
\begin{verbatim}
git status
git branch
git pull origin main
\end{verbatim}
\end{block}
\begin{exampleblock}{Meaning}
\texttt{git pull} retrieves remote changes and integrates them into the current local branch. Read the output to determine whether the branch was already current, fast-forwarded, or requires conflict resolution.
\end{exampleblock}
\begin{alertblock}{Before Pulling}
Start from a clean, understood working state. If local edits are unfinished, commit them appropriately or follow the instructor's recovery procedure rather than forcing the pull.
\end{alertblock}
\end{frame}

\begin{frame}[fragile]{Publish Local Commits with \texttt{git push}}
\framesubtitle{Synchronize verified history to GitHub}
\footnotesize
\begin{block}{Purpose and Syntax}
\begin{verbatim}
git push <remote> <branch>
\end{verbatim}
\end{block}
\begin{exampleblock}{Examples}
\begin{verbatim}
git push origin main
# For a new contribution branch:
git push -u origin docs/student-git-notes
\end{verbatim}
\end{exampleblock}
\begin{alertblock}{Rejected Push}
A non-fast-forward rejection usually means the remote contains commits missing locally. Do not force-push. Retrieve and understand upstream changes, resolve any conflict, test, and then push normally.
\end{alertblock}
\end{frame}

\begin{frame}{The Everyday Internship Workflow}
\framesubtitle{Observe, synchronize, modify, review, record, and share}
\scriptsize
\begin{enumerate}\setlength{\itemsep}{0.28em}
\item \textbf{Enter and inspect:} move to the repository, then run \texttt{git status} and \texttt{git branch}.
\item \textbf{Synchronize:} pull the approved branch before beginning new work.
\item \textbf{Modify:} make one focused code, configuration, test, or documentation change.
\item \textbf{Review:} inspect \texttt{git status} and \texttt{git diff}.
\item \textbf{Stage:} add only the files intended for the next checkpoint.
\item \textbf{Verify staged content:} inspect \texttt{git diff --staged}.
\item \textbf{Commit:} use a concise message that explains the result.
\item \textbf{Synchronize again:} pull if necessary, resolve carefully, test, and push.
\item \textbf{Confirm on GitHub:} verify the branch, commit, and documentation online.
\end{enumerate}
\begin{alertblock}{Team Rule}
The exact branch and review policy is defined by the internship repository. Follow the documented collaboration workflow rather than assuming every contributor pushes directly to \texttt{main}.
\end{alertblock}
\end{frame}

\begin{frame}[fragile]{Use \texttt{.gitignore} Deliberately}
\framesubtitle{Keep generated, local, sensitive, and oversized artifacts out of normal history}
\footnotesize
\begin{block}{Purpose}
A \texttt{.gitignore} file contains patterns for untracked files Git should normally ignore.
\end{block}
\begin{exampleblock}{Typical Robotics-Oriented Patterns}
\begin{verbatim}
# Python cache
__pycache__/
*.pyc

# Local environments and generated outputs
.venv/
build/
install/
log/

# Local secrets
.env
*.pem
\end{verbatim}
\end{exampleblock}
\begin{alertblock}{Important Boundary}
Ignoring a file does not remove it from history if it was already committed. Never rely on \texttt{.gitignore} as the only protection for credentials.
\end{alertblock}
\end{frame}

\begin{frame}{Never Commit Secrets}
\framesubtitle{Repository history can outlive the working file}
\footnotesize
\begin{alertblock}{Do Not Commit}
\begin{itemize}\setlength{\itemsep}{0.20em}
\item passwords, access tokens, API keys, or private SSH keys;
\item private certificates or machine-specific credentials;
\item confidential datasets or personal information;
\item environment files containing secret values.
\end{itemize}
\end{alertblock}
\begin{block}{If a Secret Is Exposed}
Stop sharing it, notify the responsible instructor or repository maintainer, revoke or rotate the credential, and follow the approved history-cleaning procedure. Merely deleting the file in a later commit may leave the secret in earlier history.
\end{block}
\begin{exampleblock}{Better Pattern}
Commit a documented template such as \texttt{.env.example} with placeholder names, then keep real secret values outside version control.
\end{exampleblock}
\end{frame}

\begin{frame}{Keep the Repository Engineering-Friendly}
\framesubtitle{Version the source of truth, not every generated artifact}
\footnotesize
\begin{columns}[T,onlytextwidth]
\begin{column}{0.48\textwidth}
\begin{block}{Good Candidates}
\begin{itemize}\setlength{\itemsep}{0.18em}
\item source code and launch files;
\item human-maintained configuration;
\item small test fixtures;
\item README files and reports;
\item scripts needed to reproduce outputs.
\end{itemize}
\end{block}
\end{column}
\begin{column}{0.48\textwidth}
\begin{exampleblock}{Review Before Adding}
\begin{itemize}\setlength{\itemsep}{0.18em}
\item build and cache directories;
\item large videos, models, bags, or datasets;
\item temporary logs and screenshots;
\item machine-specific settings;
\item files obtainable from documented sources.
\end{itemize}
\end{exampleblock}
\end{column}
\end{columns}
\begin{alertblock}{Reason}
Large binary files make cloning, reviewing, and history management harder. Use the repository's approved storage and release strategy for necessary large assets.
\end{alertblock}
\end{frame}

\begin{frame}[fragile]{Recover: Staged the Wrong File}
\framesubtitle{Correct the selection without destroying the working change}
\footnotesize
\begin{exampleblock}{Situation}
You staged a file that should not be part of the next commit.
\begin{verbatim}
git status
git restore --staged path/to/file
git status
\end{verbatim}
\end{exampleblock}
\begin{block}{Effect}
The file is removed from the staging area, but its working-tree modification is preserved for later review.
\end{block}
\begin{alertblock}{Do Not Guess Recovery Commands}
Always inspect \texttt{git status} before and after recovery. Commands involving \texttt{--hard}, forced pushes, or history rewriting are outside the beginner workflow unless supervised.
\end{alertblock}
\end{frame}

\begin{frame}[fragile]{Recover: Discard an Unwanted Local Edit}
\framesubtitle{Use only when the last committed version is truly the desired state}
\footnotesize
\begin{exampleblock}{For One Tracked, Unstaged File}
\begin{verbatim}
git diff -- path/to/file
git restore path/to/file
git status
\end{verbatim}
\end{exampleblock}
\begin{alertblock}{Irreversible Working-Tree Loss}
\texttt{git restore path/to/file} discards the selected uncommitted edit. Copy valuable experimental content elsewhere or commit it on an appropriate branch before restoring.
\end{alertblock}
\begin{block}{Safer Question}
Do I want to remove the file from the next commit, or destroy the actual edit? These are different intentions and require different commands.
\end{block}
\end{frame}

\begin{frame}[fragile]{Recover: Push Rejected or Pull Conflicted}
\framesubtitle{Preserve evidence and resolve the real synchronization problem}
\scriptsize
\begin{columns}[T,onlytextwidth]
\begin{column}{0.48\textwidth}
\begin{block}{Push Rejected}
\begin{verbatim}
git status
git branch
git remote -v
git pull origin main
\end{verbatim}
Read the output, integrate approved remote work, test, and push again. Never use force as the automatic response.
\end{block}
\end{column}
\begin{column}{0.48\textwidth}
\begin{exampleblock}{Merge Conflict}
\begin{enumerate}\setlength{\itemsep}{0.15em}
\item run \texttt{git status};
\item open each conflicted file;
\item understand both alternatives;
\item remove conflict markers after deciding;
\item test the integrated result;
\item stage and commit the resolution.
\end{enumerate}
\end{exampleblock}
\end{column}
\end{columns}
\begin{alertblock}{Escalation Rule}
If the correct engineering intention is unclear, stop and ask the file owner or instructor. A syntactically resolved conflict can still create incorrect robot behavior.
\end{alertblock}
\end{frame}

\begin{frame}[fragile]{Two Computers, One Shared Repository}
\framesubtitle{Synchronize through commits, not by copying folders manually}
\footnotesize
\begin{block}{At the Computer Where Work Was Completed}
\begin{verbatim}
git status
git add <intended-files>
git commit -m "describe the completed change"
git pull origin main
git push origin main
\end{verbatim}
\end{block}
\begin{exampleblock}{At the Other Computer}
\begin{verbatim}
cd ~/quard_um6p_intership
git status
git pull origin main
\end{verbatim}
\end{exampleblock}
\begin{alertblock}{Precondition}
Each computer must be on the intended repository and branch. Uncommitted local edits must be reviewed before pulling, or they may conflict with the incoming work.
\end{alertblock}
\end{frame}

\begin{frame}[fragile]{End-to-End Student Practice}
\framesubtitle{Obtain, modify, record, synchronize, and verify}
\scriptsize
\begin{exampleblock}{Guided Workflow}
\begin{verbatim}
cd ~
git clone https://github.com/abdtnourji/quard_um6p_intership.git
cd quard_um6p_intership

git status
git remote -v
git log --oneline -5

mkdir -p student_work
echo "Git practice completed" > student_work/git_practice.txt

git status
git add student_work/git_practice.txt
git diff --staged
git commit -m "docs: add Git practice evidence"
git pull origin main
git push origin main
\end{verbatim}
\end{exampleblock}
\begin{alertblock}{Access Boundary}
If direct push access or pushing to \texttt{main} is not permitted, create the approved branch and follow the repository's contribution or pull-request workflow.
\end{alertblock}
\end{frame}

\begin{frame}{Git and GitHub Readiness Check}
\framesubtitle{Can the student explain both the commands and the system state?}
\footnotesize
\begin{columns}[T,onlytextwidth]
\begin{column}{0.48\textwidth}
\begin{block}{Conceptual Understanding}
\begin{enumerate}\setlength{\itemsep}{0.18em}
\item How are Git and GitHub different?
\item What is the staging area?
\item Why is a commit local before a push?
\item What does \texttt{origin} represent?
\item Why can a merge require judgment?
\end{enumerate}
\end{block}
\end{column}
\begin{column}{0.48\textwidth}
\begin{exampleblock}{Practical Capability}
\begin{enumerate}\setlength{\itemsep}{0.18em}
\item clone and verify the repository;
\item read its documentation and structure;
\item inspect, stage, and commit a change;
\item pull and push safely;
\item recover from simple staging or editing mistakes.
\end{enumerate}
\end{exampleblock}
\end{column}
\end{columns}
\end{frame}

\begin{frame}{Good Practices for the Internship}
\framesubtitle{A concise collaboration contract}
\footnotesize
\begin{itemize}\setlength{\itemsep}{0.28em}
\item begin with \texttt{git status}, the correct branch, and an updated local repository;
\item make focused changes and inspect diffs before staging;
\item commit coherent work with meaningful messages;
\item synchronize frequently enough to avoid long divergence;
\item update documentation whenever setup or behavior changes;
\item keep generated and unnecessary large files out of normal history;
\item never commit passwords, tokens, credentials, or private keys;
\item test integrated robotics behavior after merges and conflict resolution;
\item verify the published result on GitHub.
\end{itemize}
\begin{alertblock}{Definition of Done}
A contribution is complete when another engineer can understand the change, reproduce the result, and identify the evidence supporting it.
\end{alertblock}
\end{frame}

\begin{frame}{Official References}
\framesubtitle{Primary documentation for continued practice}
\scriptsize
\begin{thebibliography}{9}
\bibitem{gitdocs}
Git Project, \emph{Git Documentation}.\\
\url{https://git-scm.com/doc}

\bibitem{gittutorial}
Git Project, \emph{A Tutorial Introduction to Git}.\\
\url{https://git-scm.com/docs/gittutorial}

\bibitem{githubbasics}
GitHub Docs, \emph{Git Basics}.\\
\url{https://docs.github.com/en/get-started/git-basics}

\bibitem{githubclone}
GitHub Docs, \emph{Cloning a Repository}.\\
\url{https://docs.github.com/en/repositories/creating-and-managing-repositories/cloning-a-repository}

\bibitem{githubpush}
GitHub Docs, \emph{Pushing Commits to a Remote Repository}.\\
\url{https://docs.github.com/en/get-started/using-git/pushing-commits-to-a-remote-repository}
\end{thebibliography}
\end{frame}

\begin{frame}{Transition to Python for Robotics}
\framesubtitle{The project is reproducible; next we study the language used for robotics logic}
\footnotesize
\begin{center}
\colorbox{um6p-color-dark-green!18}{\parbox[c][0.95cm][c]{2.05cm}{\centering\textbf{Ubuntu}\\\scriptsize Completed}}
\quad$\Longrightarrow$\quad
\colorbox{um6p-color-dark-green!18}{\parbox[c][0.95cm][c]{2.15cm}{\centering\textbf{Terminal and Bash}\\\scriptsize Completed}}
\quad$\Longrightarrow$\quad
\colorbox{um6p-color-dark-green!18}{\parbox[c][0.95cm][c]{2.05cm}{\centering\textbf{Git and GitHub}\\\scriptsize Completed}}
\end{center}
\vspace{0.18cm}
\begin{center}
$\Downarrow$

\colorbox{um6p-color-light-blue}{\parbox[c][1.08cm][c]{2.45cm}{\centering\textbf{Python for Robotics}\\\scriptsize Next foundation}}
\end{center}
\begin{exampleblock}{Guiding Question}
How can Python express readable mission logic, process robotics data, implement ROS 2 nodes, and connect camera perception to the wider inspection system?
\end{exampleblock}
\end{frame}














\begin{frame}{Transition to the Technical Foundations}
\framesubtitle{One foundation completed, the next one begins}

% \footnotesize

% ------------------------------------------------------------------
% CURRENT PROGRESS
% ------------------------------------------------------------------
\begin{center}
    \textbf{\color{um6p-dark}Internship Learning Roadmap}
\end{center}

\vspace{0.15cm}

\begin{center}
\renewcommand{\arraystretch}{1.15}
\begin{tabular}{c@{\hspace{0.10cm}}c@{\hspace{0.10cm}}c@{\hspace{0.10cm}}c@{\hspace{0.10cm}}c}

% Completed foundation
\colorbox{um6p-color-dark-green!20}{
    \parbox[c][0.85cm][c]{1.65cm}{
        \centering
        \scriptsize\bfseries
        Ubuntu\\
        \textcolor{um6p-color-dark-green}{Completed}
    }
}


& $\longrightarrow$ &

% Completed foundation
\colorbox{um6p-color-dark-green!20}{
    \parbox[c][1.05cm][c]{2.15cm}{
        \centering
        \scriptsize\bfseries
        Terminal\\
        and Bash\\[-0.05cm]
        \textcolor{um6p-color-dark-green}{Completed}
    }
}

& $\longrightarrow$ &

% Upcoming foundations
\colorbox{um6p-color-dark-green!20}{
    \parbox[c][0.85cm][c]{1.65cm}{
        \centering
        \scriptsize
        Git and\\
        GitHub \\[-0.05cm]
        \textcolor{um6p-color-dark-green}{Completed}
    }
}

\\[0.25cm]

\colorbox{um6p-color-light-blue}{
    \parbox[c][0.85cm][c]{1.65cm}{
        \centering
        \scriptsize
        Python for\\
        Robotics \\ [-0.05cm]
        \textcolor{um6p-dark}{Current}

    }
}

& $\longrightarrow$ &

\colorbox{black!7}{
    \parbox[c][0.85cm][c]{2.15cm}{
        \centering
        \scriptsize
        ROS 2
    }
}

& $\longrightarrow$ &

\colorbox{black!7}{
    \parbox[c][0.85cm][c]{1.65cm}{
        \centering
        \scriptsize
        Gazebo\\
        and PX4
    }
}

\end{tabular}

\vspace{0.15cm}

\colorbox{black!7}{
    \parbox[c][0.75cm][c]{2.1cm}{
        \centering
        \scriptsize
        AI Integration
    }
}
\end{center}


\end{frame}
